\documentclass[10pt]{article}
\usepackage[margin=1in]{geometry}
\usepackage{times}
\usepackage{graphicx}
\usepackage{amsmath}
\usepackage{amssymb}
\usepackage{hyperref}
\usepackage{enumitem}
\usepackage{booktabs}
\usepackage{tikz}
\usepackage{xcolor}
\usepackage{titlesec}
\usepackage{fancyhdr}
\usepackage{array}

% Header
\pagestyle{fancy}
\fancyhf{}
\rhead{Mobile Computing and Its Applications}
\lhead{Mini Project Proposal}
\cfoot{\thepage}
\renewcommand{\headrulewidth}{0.4pt}

% Section formatting
\titleformat{\section}{\normalfont\large\bfseries}{\thesection}{1em}{}
\titleformat{\subsection}{\normalfont\normalsize\bfseries}{\thesubsection}{1em}{}

\begin{document}

% -----------------------------------------------
%  Title Block
% -----------------------------------------------
\begin{center}
  {\LARGE \textbf{Lie \& Emotion Detector}} \\[6pt]
  {\large Real-time Deception Detection via Facial Expression Analysis} \\[10pt]
  {\normalsize Mobile Computing and Its Applications (4190.406B) --- Spring 2026} \\
  {\normalsize Prof. Youngki Lee \quad Seoul National University} \\[6pt]
  \rule{\linewidth}{0.4pt}
\end{center}
\vspace{6pt}

% -----------------------------------------------
%  1. Goal & Usage Scenario
% -----------------------------------------------
\section{Goal \& Usage Scenario}

The goal of this project is to build a real-time mobile app that detects whether a user's
displayed emotion is genuine, and if not, infers the underlying true emotion using camera
input alone.

\textbf{Final output:}
\begin{itemize}[leftmargin=*, topsep=2pt, itemsep=1pt]
  \item Is this emotion \textbf{genuine or deceptive}? (True / False)
  \item If deceptive, what is the \textbf{true emotion}? (e.g., angry, sad, \ldots)
\end{itemize}

\textbf{Usage scenarios:} interview practice, negotiation, social skills training, or
simply checking whether your friend is really fine.

% -----------------------------------------------
%  2. Related Works
% -----------------------------------------------
\section{Related Works}

\textbf{Micro-expression recognition.}
CASME II~\cite{yan2014casme2} and SAMM~\cite{davison2018samm} are widely used benchmarks
for spontaneous micro-expression detection, where subtle involuntary expressions reveal
concealed emotions.

\textbf{Deception detection.}
The MU3D dataset~\cite{lloyd2019mu3d} contains videos of genuine and deceptive facial
expressions across multiple emotion categories, making it well suited as a
supervised signal for deception classification.

\textbf{Semi-supervised learning.}
FixMatch~\cite{sohn2020fixmatch} achieves strong performance with few labeled examples
by assigning pseudo-labels to unlabeled data based on a confidence threshold.
SimCLR~\cite{chen2020simclr} learns transferable visual representations via
contrastive learning without labels.

\textbf{Lightweight on-device inference.}
MediaPipe Face Mesh~\cite{lugaresi2019mediapipe} extracts 468 3-D facial landmarks
in real time on mobile devices, enabling lightweight coordinate-based DNN inputs
instead of raw image inputs.

% -----------------------------------------------
%  3. Key Idea
% -----------------------------------------------
\section{Key Idea}

\subsection{Pipeline Overview}

The system is composed of three models, summarized as $A + B \rightarrow C$:

\begin{center}
\small
\begin{tabular}{cl}
\toprule
\textbf{Model} & \textbf{Role} \\
\midrule
$A$ & Facial expression $\rightarrow$ Emotion (7 classes) \\
$B$ & Genuine vs.\ deceptive classification \\
$C$ & Fusion: true emotion given deception signal \\
\bottomrule
\end{tabular}
\end{center}

\vspace{4pt}
MediaPipe Face Mesh extracts 468 landmark coordinates per frame. These coordinates
(not raw pixels) serve as inputs to Models $A$ and $B$, making both models extremely
lightweight and quantization-friendly.

\subsection{Semi-Supervised Learning with FixMatch}

A key challenge is the limited availability of labeled deception data.
We address this with the following four-stage pipeline:

\begin{enumerate}[leftmargin=*, topsep=2pt, itemsep=2pt]
  \item \textbf{Contrastive pre-training (SimCLR).}
        Train a feature extractor on \emph{both} labeled (Dataset $B$: self-collected)
        and unlabeled data (Dataset $A$: MU3D used without labels) via
        contrastive loss. This produces a shared facial-feature embedding space.

  \item \textbf{Supervised fine-tuning.}
        Attach a classifier head and train only on Dataset $B$ (labeled),
        using the frozen SimCLR backbone.

  \item \textbf{Pseudo-labeling (FixMatch).}
        Apply the fine-tuned model to Dataset $A$. Retain pseudo-labels only
        where prediction confidence exceeds threshold $\tau = 0.95$.

  \item \textbf{Joint re-training.}
        Combine pseudo-labeled $A$ with labeled $B$ and retrain the full
        Model $B$ end-to-end.
\end{enumerate}

This allows us to exploit the large unlabeled MU3D corpus to improve deception
detection without requiring manual annotation.

% -----------------------------------------------
%  4. Implementation Plan
% -----------------------------------------------
\section{Implementation Plan}

\subsection{Datasets}

\begin{itemize}[leftmargin=*, topsep=2pt, itemsep=1pt]
  \item \textbf{AffectNet}: 450k images, 7 emotion classes $\rightarrow$ Model $A$
  \item \textbf{MU3D}: genuine/deceptive expression videos $\rightarrow$ unlabeled pool for FixMatch
  \item \textbf{Self-collected}: minimum 20 subjects; genuine vs.\ posed smiles (5 s clips),
        controlled lighting/distance $\rightarrow$ labeled set for Model $B$
\end{itemize}

\subsection{Model Architecture}

All three models use \textbf{MobileNetV2} backbones fine-tuned on landmark coordinate
sequences. Model $C$ is a lightweight MLP that takes the concatenated output logits
of $A$ and $B$ as input.

\subsection{Model Optimization}

Models will be converted to TFLite with \textbf{INT8 post-training quantization}.
Because inputs are low-dimensional landmark coordinates rather than images,
accuracy degradation is expected to be $<2\%$ while achieving $\sim75\%$
size reduction --- maximizing the scoring criteria for size reduction rate vs.\
accuracy drop.

\subsection{Android App}

\begin{itemize}[leftmargin=*, topsep=2pt, itemsep=1pt]
  \item Real-time camera feed with MediaPipe Face Mesh overlay
  \item Live emotion gauge (surface emotion vs.\ inferred true emotion)
  \item Deception probability bar; red border alert on high probability
  \item Post-session summary heatmap of expression timeline
\end{itemize}

% -----------------------------------------------
%  5. Project Timeline
% -----------------------------------------------
\section{Project Timeline}

\begin{center}
\small
\begin{tabular}{>{\bfseries}p{2.0cm} p{4.8cm}}
\toprule
Period & Task \\
\midrule
$\sim$May 4   & Proposal submission; environment setup; dataset download \\
May 5--11      & SimCLR pre-training; Model A \& B initial training \\
May 12--18     & Self-data collection; FixMatch pseudo-labeling; Model C fusion \\
May 19--25     & TFLite INT8 quantization; accuracy/size benchmarking \\
May 26--Jun 1  & Android app + visualization; GitHub cleanup; final APK \\
Jun 2--4       & Final presentation \\
\bottomrule
\end{tabular}
\end{center}

\subsection*{Team Role Division}

\begin{center}
\small
\begin{tabular}{p{3.5cm} p{3.5cm}}
\toprule
\textbf{Member A} & \textbf{Member B} \\
\midrule
Model training (Python) & Android app development \\
SimCLR + FixMatch & MediaPipe integration \\
TFLite quantization & UI \& visualization \\
\multicolumn{2}{c}{Self-data collection (joint)} \\
\bottomrule
\end{tabular}
\end{center}

% -----------------------------------------------
%  References
% -----------------------------------------------
\bibliographystyle{plain}
\begin{thebibliography}{9}

\bibitem{yan2014casme2}
W.~Yan et al.,
``CASME II: An Improved Spontaneous Micro-Expression Database,''
\textit{PLOS ONE}, 2014.

\bibitem{davison2018samm}
A.~K. Davison et al.,
``SAMM: A Spontaneous Micro-Facial Movement Dataset,''
\textit{IEEE T-AFFC}, 2018.

\bibitem{lloyd2019mu3d}
E.~P. Lloyd et al.,
``The MU3D: A Multi-Emotion, Multi-Intensity, 3D Face Database,''
\textit{PLOS ONE}, 2019.

\bibitem{sohn2020fixmatch}
K.~Sohn et al.,
``FixMatch: Simplifying Semi-Supervised Learning with Consistency and Confidence,''
\textit{NeurIPS}, 2020.

\bibitem{chen2020simclr}
T.~Chen et al.,
``A Simple Framework for Contrastive Learning of Visual Representations,''
\textit{ICML}, 2020.

\bibitem{lugaresi2019mediapipe}
C.~Lugaresi et al.,
``MediaPipe: A Framework for Building Perception Pipelines,''
\textit{arXiv:1906.08172}, 2019.

\end{thebibliography}

\end{document}
